\section{CRT factorisation at the transfer-matrix level}
\label{sec:crt_factor}

This section establishes the tensor decomposition of the transfer
matrix $T_m$ and of its Ruelle measure
$\pi_m^{\mathrm{Ruelle}}$. The decomposition was anticipated in
Section~\ref{ssec:transfer_matrix}; here we prove it and verify it
numerically.

\subsection{The Chinese Remainder Theorem as a ring isomorphism}
\label{ssec:crt_ring}

For any squarefree integer $m = p_1 \cdots p_k$ with distinct
primes $p_i$, the map
\begin{equation}
  \chi_m \,:\;
  \mathbb{Z}/m\mathbb{Z}
  \;\longrightarrow\;
  \prod_{i=1}^{k} \mathbb{Z}/p_i\mathbb{Z},
  \qquad
  r \;\longmapsto\; (r \bmod p_1, \ldots, r \bmod p_k)
  \label{eq:CRT_iso}
\end{equation}
is a ring isomorphism~\cite[Thm.~2.18]{Davenport2000}. We will use
$\chi_m$ without comment to identify residues mod~$m$ with $k$-tuples
of residues mod~$p_i$, and conversely.

The identification~\eqref{eq:CRT_iso} restricts to an identification
of the active subspaces:
$\chi_m\bigl(\mathcal{S}_m\bigr) = \prod_i \mathcal{S}_{p_i}$, where
$\mathcal{S}_p = \{1, \ldots, p-1\}$ as in
Section~\ref{ssec:transfer_matrix}. This is the structural reason
behind the tensor decomposition that follows: an active residue
class mod~$m$ is, by CRT, a tuple of active residue classes mod the
prime factors.

\subsection{The tensor decomposition of $T_m$}
\label{ssec:Tm_decomposition}

\begin{theorem}[CRT Tensor Factorization]
\label{thm:crt_tensor}
\footnote{Formally verified in Lean~4 + Mathlib as
\texttt{PT.CrtDecoupling.Tensor.crt\_tensor\_factor};
see Section~\ref{ssec:formal_verification}.}
For $m = p_1 \cdots p_k$ a product of distinct primes, the
transfer matrix factorises tensorially:
\begin{equation}
  T_m \;=\; \bigotimes_{i=1}^{k} T_{p_i}.
  \label{eq:Tm_tensor}
\end{equation}
\end{theorem}

\begin{proof}
Identify $\mathcal{S}_m$ with $\prod_i \mathcal{S}_{p_i}$ via
$\chi_m$, and write $\mathbf{r} = (r^{(1)}, \ldots, r^{(k)})$,
$\mathbf{r}' = (r^{(1)\prime}, \ldots, r^{(k)\prime})$ for
elements of the product. The transition kernel of the prime-sieve
Markov chain on $\mathcal{S}_m$ is the joint kernel of $k$
independent prime-sieve chains on the factor spaces: for
$g \sim \mathrm{Geom}(q)$ a geometric prime gap drawn from the
sieve dynamics,
\begin{equation}
  \Pr\!\bigl[\mathbf{r}' \mid \mathbf{r}\bigr]
  \;=\;
  \Pr\!\bigl[g \equiv r^{(i)\prime} - r^{(i)} \pmod{p_i}
        \text{ for all } i\bigr].
\end{equation}
By CRT, the simultaneous congruence
$g \equiv r^{(i)\prime} - r^{(i)} \pmod{p_i}$ for $i = 1, \ldots, k$
is equivalent to a single congruence
$g \equiv \chi_m^{-1}(\mathbf{r}' - \mathbf{r}) \pmod{m}$, so the
joint kernel splits into a product
$\prod_i (T_{p_i})_{r^{(i)}, r^{(i)\prime}}$. This is the matrix
entry of the Kronecker product $\bigotimes_i T_{p_i}$ at the
indices $(\mathbf{r}, \mathbf{r}')$, which proves
Eq.~\eqref{eq:Tm_tensor}.
\end{proof}

The factorisation is well known in the persistence-theory
literature; see in particular~\cite[eq.~CRT\_factor]{PT_GRAVITON}
and~\cite[Lemma~G, step~G.2]{PT_Monograph}. We give the short proof
here for self-containedness and to fix the active-subspace
convention used throughout the paper.

\subsection{The tensor decomposition of $\pi_m^{\mathrm{Ruelle}}$}
\label{ssec:pi_decomposition}

\begin{lemma}[Ruelle Factorization]
\label{lem:pi_factor}
\footnote{Formally verified in Lean~4 + Mathlib as
\texttt{PT.CrtDecoupling.Tensor.ruelle\_factor};
see Section~\ref{ssec:formal_verification}.}
Let $T_p$ have unique left eigenvector $\pi_p^{\mathrm{Ruelle}}$
for its Perron eigenvalue $\lambda_0 = 1$, normalised to a
probability distribution. Then
\begin{equation}
  \pi_{m}^{\mathrm{Ruelle}}
  \;=\;
  \bigotimes_{i=1}^{k} \pi_{p_i}^{\mathrm{Ruelle}}.
  \label{eq:pi_factor}
\end{equation}
\end{lemma}

\begin{proof}
By the Kronecker product rule,
$(\pi_p \otimes \pi_q)(T_p \otimes T_q) = (\pi_p T_p) \otimes
(\pi_q T_q) = \pi_p \otimes \pi_q$, so the right-hand side
of~\eqref{eq:pi_factor} is a left eigenvector of $T_m$ for
$\lambda = 1$.
The spectrum of $T_p \otimes T_q$ is
$\{\lambda_i^{(p)} \lambda_j^{(q)}\}_{i,j}$, and the value~$1$
is attained uniquely by $(\lambda_0^{(p)}, \lambda_0^{(q)}) =
(1, 1)$ since each $T_p$ has spectral gap (T4 corpus,
\cite[Ch.~7]{PT_Monograph}).
Uniqueness then follows from Perron--Frobenius primitivity, which
is preserved by tensor products (Appendix~\ref{appssec:tensor_primitive}).
\end{proof}

For $m = pq$ with $p, q \in \{3, 5, 7\}$ distinct, the
factorisation~\eqref{eq:pi_factor} reads
$\pi_{pq}^{\mathrm{Ruelle}} = \pi_p^{\mathrm{Ruelle}} \otimes
\pi_q^{\mathrm{Ruelle}}$ with each marginal uniform on the
respective active subspace; this is the form used in the proof of
Theorem~\ref{thm:ruelle_zero} in Section~\ref{ssec:thm_4_2_proof}.

\subsection{Numerical verification}
\label{ssec:t2_verification}

The factorisation~\eqref{eq:pi_factor} is verified numerically by
the test \texttt{T2} of
script \texttt{scripts/test\_crt\_decoupling.py}. The script
\begin{enumerate}[itemsep=2pt,leftmargin=*]
  \item builds $T_3 = \mathrm{antidiag}(1, 1)$ on
    $\mathcal{S}_3 = \{1, 2\}$ and $T_5$ on
    $\mathcal{S}_5 = \{1, 2, 3, 4\}$ with off-diagonal entries
    $1/(p-2)$ as in Eq.~\eqref{eq:Tp_explicit};
  \item solves the left eigenvector equation $\pi_p T_p = \pi_p$
    on each factor by augmented least squares
    at \texttt{mpmath} precision $50$~digits;
  \item builds $T_{15} = T_3 \otimes T_5$ as an $8 \times 8$
    Kronecker product, and likewise builds the candidate
    factorised eigenvector $\pi_3 \otimes \pi_5$;
  \item independently solves the left eigenvector equation
    $\pi_{15} T_{15} = \pi_{15}$ at the same precision;
  \item reports the $L^2$ norm of the residual
    $\pi_{15}^{\mathrm{direct}} - \pi_3 \otimes \pi_5$ and the
    KL divergence
    $\DKL(\pi_{15}^{\mathrm{direct}} \| \pi_3 \otimes \pi_5)$.
\end{enumerate}
The observed residual is
\begin{equation}
  \bigl\|\,\pi_{15}^{\mathrm{direct}}
      - \pi_3^{\mathrm{Ruelle}} \otimes \pi_5^{\mathrm{Ruelle}}\,\bigr\|_2
  \;=\; 5.88 \times 10^{-53}
  \qquad \text{at 50 dps},
\end{equation}
and the KL divergence is zero to the same precision. We record this
as the numerical witness of
Lemma~\ref{lem:pi_factor}; see Appendix~\ref{appssec:T2} for the
full script output.
